% GENERATED FILE. Edit canonical lessons, then run npm run generate:books.
% canonical-source-hash: fnv1a64:e4a81d87c7a44de2
% canonical-lessons: GE-C249-ergaenzen, GE-C249-erledigen, GE-C249-erscheinen, GE-C249-erzeugen, GE-C249-fordern

\chapter{To complete, To get done, To appear, To produce, To demand}
\label{ch:ge-to-complete-to-get-done-to-appear-to-produce-to-demand}
\hlchaptermodality{249}

\begin{chapteropening}
\textbf{By the end of this chapter:} \emph{I can say to complete, to get done, to appear, to produce and to demand.}

Complete the last lesson of chapter 249: i can say to complete, to get done, to appear, to produce and to demand.
\end{chapteropening}

\section[ergänzen]{ergänzen — to complete, to add}

\label{lesson:GE-C249-ergaenzen}



\begin{quote}
Before the new one: say the German for to excuse, then the German for to find out.
\end{quote}

\subsection*{You'll want to know: ergänzen}
\textbf{ergänzen} — \textquotedblleft{}to complete, to add\textquotedblright{}.

\begin{grammarlens}[title={What you've built}]
One more word for this chapter.
\end{grammarlens}

\subsection*{Your turn}
\begin{itemize}
\raggedright
  \item \emph{Say it:} \emph{ergänzen}
  \item \emph{Say it:} \emph{ergänzen}, once more
  \item \emph{Say it:} say \emph{ergänzen}
  \item \emph{Recall:} say the German for to excuse, then the German for to find out, then say \emph{ergänzen} again
\end{itemize}

\begin{tcolorbox}[breakable,colback=teal!4,colframe=teal!35!black,title={Before you move on}]
What does ergänzen mean? (\textquotedblleft{}To complete, to add\textquotedblright{}.) Say it once more.
\end{tcolorbox}
\section[erledigen]{erledigen — to get done}

\label{lesson:GE-C249-erledigen}



\begin{quote}
Before the new one: say the German for to find out, then the German for to complete.
\end{quote}

\subsection*{You'll want to know: erledigen}
\textbf{erledigen} — \textquotedblleft{}to get done\textquotedblright{}.

\begin{grammarlens}[title={What you've built}]
One more word for this chapter.
\end{grammarlens}

\subsection*{Your turn}
\begin{itemize}
\raggedright
  \item \emph{Say it:} \emph{erledigen}
  \item \emph{Say it:} \emph{erledigen}, once more
  \item \emph{Say it:} say \emph{erledigen}
  \item \emph{Recall:} say the German for to find out, then the German for to complete, then say \emph{erledigen} again
\end{itemize}

\begin{tcolorbox}[breakable,colback=teal!4,colframe=teal!35!black,title={Before you move on}]
What does erledigen mean? (\textquotedblleft{}To get done\textquotedblright{}.) Say it once more.
\end{tcolorbox}
\section[erscheinen]{erscheinen — to appear}

\label{lesson:GE-C249-erscheinen}



\begin{quote}
Before the new one: say the German for to complete, then the German for to get done.
\end{quote}

\subsection*{You'll want to know: erscheinen}
\textbf{erscheinen} — \textquotedblleft{}to appear\textquotedblright{}.

\begin{grammarlens}[title={What you've built}]
One more word for this chapter.
\end{grammarlens}

\subsection*{Your turn}
\begin{itemize}
\raggedright
  \item \emph{Say it:} \emph{erscheinen}
  \item \emph{Say it:} \emph{erscheinen}, once more
  \item \emph{Say it:} say \emph{erscheinen}
  \item \emph{Recall:} say the German for to complete, then the German for to get done, then say \emph{erscheinen} again
\end{itemize}

\begin{tcolorbox}[breakable,colback=teal!4,colframe=teal!35!black,title={Before you move on}]
What does erscheinen mean? (\textquotedblleft{}To appear\textquotedblright{}.) Say it once more.
\end{tcolorbox}
\section[erzeugen]{erzeugen — to produce}

\label{lesson:GE-C249-erzeugen}



\begin{quote}
Before the new one: say the German for to get done, then the German for to appear.
\end{quote}

\subsection*{You'll want to know: erzeugen}
\textbf{erzeugen} — \textquotedblleft{}to produce\textquotedblright{}.

\begin{grammarlens}[title={What you've built}]
One more word for this chapter.
\end{grammarlens}

\subsection*{Your turn}
\begin{itemize}
\raggedright
  \item \emph{Say it:} \emph{erzeugen}
  \item \emph{Say it:} \emph{erzeugen}, once more
  \item \emph{Say it:} say \emph{erzeugen}
  \item \emph{Recall:} say the German for to get done, then the German for to appear, then say \emph{erzeugen} again
\end{itemize}

\begin{tcolorbox}[breakable,colback=teal!4,colframe=teal!35!black,title={Before you move on}]
What does erzeugen mean? (\textquotedblleft{}To produce\textquotedblright{}.) Say it once more.
\end{tcolorbox}
\section[fordern]{fordern — to demand}

\label{lesson:GE-C249-fordern}



\begin{quote}
Before the new one: say the German for to appear, then the German for to produce.
\end{quote}

\subsection*{You'll want to know: fordern}
\textbf{fordern} — \textquotedblleft{}to demand\textquotedblright{}.

\begin{grammarlens}[title={What you've built}]
One more word for this chapter.
\end{grammarlens}

\subsection*{Your turn}
\begin{itemize}
\raggedright
  \item \emph{Say it:} \emph{fordern}
  \item \emph{Say it:} \emph{fordern}, once more
  \item \emph{Say it:} say \emph{fordern}
  \item \emph{Recall:} say the German for to appear, then the German for to produce, then say \emph{fordern} again
\end{itemize}

\begin{tcolorbox}[breakable,colback=teal!4,colframe=teal!35!black,title={Before you move on}]
What does fordern mean? (\textquotedblleft{}To demand\textquotedblright{}.) Say it once more.
\end{tcolorbox}
